% Figure: a hanging chain settles into the shape that minimises its
% gravitational potential energy, the catenary y = a cosh(x/a), contrasted
% with the parabola a naive analysis assumes. The two curves share
% endpoints and sag but differ in the middle.
\begin{figure}[htbp]
\centering
\begin{tikzpicture}[
  axis/.style={draw=engblack, -{Stealth}, thick},
  cat/.style={draw=engblue, very thick},
  para/.style={draw=engred, very thick, dashed},
  guide/.style={draw=enggray, dashed, thin},
  label/.style={font=\small},
]
% supports at x = -3 and x = +3, at height 3
% catenary y = a cosh(x/a) - c, choose a = 2.0
% at x=3: 2 cosh(1.5) = 2*2.3524 = 4.705; subtract offset so y(3)=3 -> c=1.705
% at x=0: y = 2 - 1.705 = 0.295 (the sag point)
\draw[axis] (-4,0) -- (4.4,0) node[anchor=west, label] {$x$};
\draw[axis] (0,-0.3) -- (0,4.2) node[anchor=south, label] {$y$};

% supports
\fill[engblack] (-3,3) circle (2.4pt);
\fill[engblack] (3,3) circle (2.4pt);
\node[font=\scriptsize, anchor=south east] at (-3,3.05) {support};
\node[font=\scriptsize, anchor=south west] at (3,3.05) {support};

% catenary (blue): y = 2 cosh(x/2) - 1.705
\draw[cat] plot[domain=-3:3, samples=120]
  ({\x}, {2.0*cosh(\x/2.0) - 1.705});
\node[engblue, label, anchor=north] at (-1.7, 0.95) {catenary};

% parabola (red dashed) through same three points (-3,3),(0,0.295),(3,3):
% y = 0.295 + k x^2, at x=3: 3 = 0.295 + 9k -> k = 0.3006
\draw[para] plot[domain=-3:3, samples=120]
  ({\x}, {0.295 + 0.3006*\x*\x});
\node[engred, label, anchor=south] at (1.9, 1.55) {parabola};

% sag marker at x = 0
\draw[guide] (0,0.295) -- (0,0) ;
\fill[englime] (0,0.295) circle (2.2pt);
\node[englime, font=\scriptsize, anchor=west] at (0.12, 0.30) {lowest point};

\end{tikzpicture}
\caption[The catenary as a minimum-energy shape]{A uniform flexible chain
hung between two supports settles into the curve that minimises its total
gravitational potential energy subject to its fixed length, the catenary
$y = a\cosh(x/a)$ (blue). The parabola (red dashed) through the same
supports and the same sag is the shape a chain would take only if its
weight were spread evenly along the horizontal rather than along its own
arc; the two curves agree at the endpoints and at the lowest point (green)
but the catenary hangs slightly fuller in between. The shape follows from
the calculus of variations applied to the potential-energy functional, the
same stationarity principle that governs the action in dynamics, with the
fixed-length condition entering through a Lagrange multiplier that turns
out to be the horizontal tension in the chain.}
\label{fig:vol03ch05:catenary}
\end{figure}
